% Variante de non-regression — en-tete article multi-etablissement.
% Les logos uftmp et n7 doivent rester sur une seule ligne, a hauteur commune.
\documentclass[11pt]{ocots-td}
\usepackage[lang=fr, theme=ocots, solutions=none, math=analysis,
            institution={uftmp,n7}]{ocots}
\lstset{language=[LaTeX]TeX}

\title{TD de démonstration — En-tête article}
\shorttitle{TD — En-tête}
\numero{TD~o 1}
\date{2025--2026}
\discipline{Calcul différentiel}
\promotion{ENSEEIHT / INSA — 1A}

\newcommand{\ocotsvariantnote}[1]{\begin{quote}\itshape Variante : #1\end{quote}}

\begin{document}

\maketitle

\begin{lstlisting}
\usepackage[..., institution={uftmp,n7}]{ocots}
\end{lstlisting}
\ocotsvariantnote{\texttt{institution=\{uftmp,n7\}} : deux établissements
séparés par une virgule et entourés d'accolades (le défaut est un seul,
\texttt{institution=n7}) — les deux logos sont composés côte à côte, à
hauteur commune, dans l'en-tête.}

\begin{instructions}
    Les logos de la première page sont composés en mode en-tête.
\end{instructions}

\begin{exercise}[points=2]
    Premier exercice : le titre doit rester immédiatement sous l'en-tête et
    l'exercice doit commencer plus haut que dans l'ancien rendu.
\end{exercise}

\end{document}
